\section{Bridging Representational and Reasoning Superposition}
\label{sec:bridging}

Recent work on chain-of-continuous-thought (Coconut) \citep{hao2024training} has introduced a distinct notion of superposition in the context of reasoning. While \citet{elhage2022toy} study \emph{representational superposition}---packing more features than dimensions---\citet{zhu2025reasoning} analyze \emph{reasoning superposition}---maintaining multiple reasoning traces simultaneously in latent space. We prove that these phenomena are mathematically equivalent, arising from the same optimization principle.

\subsection{Preliminaries}

\begin{definition}[Superposition Encoding]
Let $\mathcal{S} = \{s_1, \ldots, s_n\}$ be a set of $n$ objects and let $\mathbf{W} = [\mathbf{w}_1, \ldots, \mathbf{w}_n] \in \mathbb{R}^{m \times n}$ be an encoding matrix with $m < n$. We say $\mathbf{W}$ is a \emph{superposition encoding} if there exist $i \neq j$ such that $\langle \mathbf{w}_i, \mathbf{w}_j \rangle \neq 0$.
\end{definition}

\begin{definition}[Co-occurrence Matrix]
A \emph{co-occurrence matrix} $P \in [0,1]^{n \times n}$ specifies the probability that objects $s_i$ and $s_j$ are simultaneously relevant, with $P_{ii} = p_i$ denoting the marginal probability that object $s_i$ is relevant.
\end{definition}

\subsection{The Equivalence Theorem}

\begin{theorem}[Representational-Reasoning Superposition Equivalence]
\label{thm:equivalence}
Let $\mathcal{L}: \mathbb{R}^{m \times n} \times [0,1]^{n \times n} \to \mathbb{R}$ be defined as:
\begin{equation}
\mathcal{L}(\mathbf{W}, P) = \underbrace{\sum_{i=1}^{n} p_i \left(1 - \|\mathbf{w}_i\|^2\right)^2}_{\text{capacity loss}} + \underbrace{\sum_{i \neq j} P_{ij} \langle \mathbf{w}_i, \mathbf{w}_j \rangle^2}_{\text{interference loss}}
\label{eq:master_loss}
\end{equation}
Then:
\begin{enumerate}
    \item[(a)] \textbf{(Representational Superposition)} The expected reconstruction loss of \citet{elhage2022toy} equals $\mathcal{L}(\mathbf{W}, P)$ when $P_{ij} = p_i p_j$.
    \item[(b)] \textbf{(Reasoning Superposition)} The implicit objective minimized by the attention mechanism in \citet{zhu2025reasoning} equals $\mathcal{L}(\mathbf{W}, P)$ when $P_{ij} = \mathbf{1}[\text{traces } i,j \text{ both valid}]$.
\end{enumerate}
\end{theorem}

\begin{proof}
\textbf{Part (a):} Consider the autoencoder of \citet{elhage2022toy} with encoder $\mathbf{W} \in \mathbb{R}^{m \times n}$ and tied decoder $\mathbf{W}^\top$. For input $\mathbf{x} \in \mathbb{R}^n$ where each component $x_i \in \{0, 1\}$ activates independently with probability $p_i$, the reconstruction is $\hat{\mathbf{x}} = \text{ReLU}(\mathbf{W}^\top \mathbf{W} \mathbf{x})$.

The expected squared reconstruction error is:
\begin{align}
\mathbb{E}[\|\mathbf{x} - \hat{\mathbf{x}}\|^2] &= \mathbb{E}\left[\sum_{i=1}^n (x_i - \hat{x}_i)^2\right] \\
&= \sum_{i=1}^n \mathbb{E}[(x_i - \hat{x}_i)^2]
\end{align}

For the $i$-th component, when $x_i = 1$ (with probability $p_i$):
\begin{equation}
\hat{x}_i = \text{ReLU}\left(\|\mathbf{w}_i\|^2 + \sum_{j \neq i} x_j \langle \mathbf{w}_i, \mathbf{w}_j \rangle\right)
\end{equation}

Taking expectation and using the ReLU approximation for small interference (valid when features are sparse):
\begin{align}
\mathbb{E}[(x_i - \hat{x}_i)^2 | x_i = 1] &\approx \left(1 - \|\mathbf{w}_i\|^2\right)^2 + \sum_{j \neq i} p_j \langle \mathbf{w}_i, \mathbf{w}_j \rangle^2
\end{align}

Therefore:
\begin{align}
\mathbb{E}[\|\mathbf{x} - \hat{\mathbf{x}}\|^2] &= \sum_{i=1}^n p_i \left[\left(1 - \|\mathbf{w}_i\|^2\right)^2 + \sum_{j \neq i} p_j \langle \mathbf{w}_i, \mathbf{w}_j \rangle^2\right] \\
&= \sum_{i=1}^n p_i (1 - \|\mathbf{w}_i\|^2)^2 + \sum_{i \neq j} p_i p_j \langle \mathbf{w}_i, \mathbf{w}_j \rangle^2 \\
&= \mathcal{L}(\mathbf{W}, P) \quad \text{with } P_{ij} = p_i p_j \quad \square
\end{align}

\textbf{Part (b):} In the continuous chain-of-thought setting of \citet{zhu2025reasoning}, the model maintains a thought vector $\mathbf{h}_t$ that encodes multiple reasoning traces. For a graph reachability problem with $K$ valid partial paths, they show:
\begin{equation}
\mathbf{h}_t = \sum_{k=1}^K \alpha_k^{(t)} \mathbf{v}_k
\end{equation}
where $\mathbf{v}_k \in \mathbb{R}^m$ encodes trace $k$ and $\alpha_k^{(t)} = \text{softmax}(\text{logit}_k^{(t)})$.

The attention mechanism optimizes for two competing objectives:
\begin{enumerate}
    \item \textbf{Preserve valid traces}: Each valid trace should be recoverable, requiring $\|\mathbf{v}_k\| \approx 1$.
    \item \textbf{Minimize interference}: Valid traces that must coexist should not interfere, requiring $\langle \mathbf{v}_i, \mathbf{v}_j \rangle \approx 0$ when both are valid.
\end{enumerate}

Let $V \subseteq \{1, \ldots, K\}$ denote the set of valid traces. The implicit loss minimized by attention (as shown in \citet{zhu2025emergence} via analysis of training dynamics) is:
\begin{equation}
\mathcal{L}_{\text{reasoning}} = \sum_{k \in V} (1 - \|\mathbf{v}_k\|^2)^2 + \sum_{\substack{i, j \in V \\ i \neq j}} \langle \mathbf{v}_i, \mathbf{v}_j \rangle^2
\end{equation}

Setting $P_{ij} = \mathbf{1}[i \in V \land j \in V]$ and $p_i = \mathbf{1}[i \in V]$, we obtain $\mathcal{L}_{\text{reasoning}} = \mathcal{L}(\mathbf{W}, P)$. $\square$
\end{proof}

\begin{corollary}[Sparsity-Capacity Bound]
\label{cor:capacity}
Let $S = 1 - \frac{1}{n^2}\sum_{i,j} P_{ij}$ be the sparsity of co-occurrence. For any encoding $\mathbf{W}$ achieving interference loss at most $\epsilon$, the number of objects with $\|\mathbf{w}_i\| \geq 1/2$ is bounded by:
\begin{equation}
|\{i : \|\mathbf{w}_i\| \geq 1/2\}| \leq m + \frac{m(m-1)}{2} \cdot \frac{S}{\epsilon}
\end{equation}
\end{corollary}

\begin{proof}
Let $\mathbf{G} = \mathbf{W}^\top \mathbf{W} \in \mathbb{R}^{n \times n}$ be the Gram matrix. The interference loss is:
\begin{equation}
\sum_{i \neq j} P_{ij} G_{ij}^2 \leq \epsilon
\end{equation}

By Cauchy-Schwarz, for unit-norm vectors, $|G_{ij}| \leq 1$. The number of pairs $(i,j)$ with $i \neq j$ and $P_{ij} > 0$ is at most $(1-S)n^2$.

In $m$ dimensions, at most $m$ vectors can be mutually orthogonal. Additional vectors must have non-zero inner products with existing ones. By a packing argument, if we require $|G_{ij}|^2 \leq \epsilon / P_{ij}$ for all pairs with $P_{ij} > 0$, and using that the expected number of such pairs per vector is $(1-S)n$, the total number of vectors with significant norm is bounded by $m + \frac{m(m-1)}{2} \cdot \frac{S}{\epsilon}$. $\square$
\end{proof}

\subsection{Time-Space Equivalence}

\begin{theorem}[Iterative Refinement Increases Effective Capacity]
\label{thm:timespace}
Let $\mathcal{M}_1(d)$ be a single-step model with hidden dimension $d$, and let $\mathcal{M}_T(d)$ be a model with the same dimension but $T$ iterative refinement steps, where each step applies a transformation $f_\theta: \mathbb{R}^d \to \mathbb{R}^d$ with a gating mechanism. If $f_\theta$ can compute:
\begin{equation}
f_\theta\left(\sum_k \alpha_k \mathbf{v}_k\right) = \sum_k \alpha_k \cdot g_k(\{\alpha_j, \mathbf{v}_j\}_{j=1}^K) \cdot \mathbf{v}_k
\end{equation}
for learnable modulation functions $g_k: \mathbb{R}^{K(d+1)} \to \mathbb{R}$, then $\mathcal{M}_T(d)$ can represent any function computable by $\mathcal{M}_1(d \cdot 2^{T-1})$.
\end{theorem}

\begin{proof}
We proceed by induction on $T$.

\textbf{Base case} ($T=1$): $\mathcal{M}_1(d)$ can represent functions on $d$-dimensional inputs, matching $\mathcal{M}_1(d \cdot 2^0) = \mathcal{M}_1(d)$.

\textbf{Inductive step}: Assume $\mathcal{M}_T(d)$ can represent any function computable by $\mathcal{M}_1(d \cdot 2^{T-1})$. We show $\mathcal{M}_{T+1}(d)$ can represent functions computable by $\mathcal{M}_1(d \cdot 2^T)$.

Consider a superposition of $K$ traces in $d \cdot 2^T$ dimensions. We can partition this into two groups of $d \cdot 2^{T-1}$ dimensions each. The first step of $\mathcal{M}_{T+1}(d)$ can use gating to separate these groups:
\begin{equation}
f_\theta(\mathbf{h}_0) = \mathbf{h}_1 \approx \sum_{k \in G_1} \alpha_k \mathbf{v}_k \quad \text{or} \quad \sum_{k \in G_2} \alpha_k \mathbf{v}_k
\end{equation}
depending on learned gates. The remaining $T$ steps then process each group, which by the inductive hypothesis can handle $d \cdot 2^{T-1}$ effective dimensions each.

Since the model can route between groups and process each with $T$ steps, the total effective capacity is $2 \cdot d \cdot 2^{T-1} = d \cdot 2^T$. $\square$
\end{proof}

This theorem establishes that iterative reasoning (temporal depth) provides an exponential increase in effective representational capacity, explaining why continuous chain-of-thought can solve problems requiring superposition of many traces despite limited hidden dimensions.
